% Part: sets-functions-relations
% Chapter: induction
% Section: induction

\documentclass[../../include/open-logic-section]{subfiles}

\begin{document}

\olfileid{sfr}{ind}{int}

\olsection{పరిచయం}

\begin{explain}
ఆగమనం అనేది గణిత తర్కంలో సాధారణంగా ఉపయోగించే నిరూపణ పద్ధతి. సహజ సంఖ్యలపై ఆగమనం చేసే గణితం నుంచి ఇది మీకు తెలిసి ఉండవచ్చు. వాస్తవానికి, ఆగమనానికి అనుకూలించే గణిత, తార్కిక వస్తు సమితులు అనేక రకాలుగా ఉంటాయి.

సాధారణంగా ఆగమనం ఒక సార్వత్రిక వాదనను నిరూపించడానికి ఉపయోగపడుతుంది; అంటే ఒక రకానికి చెందిన ప్రతి వస్తువుకు ఏదో లక్షణం ఉందని చూపడానికి. ముఖ్యంగా, “సార్వత్రిక ప్రవేశం” అనే నిరూపణ పద్ధతి చాలా సంక్లిష్టమయ్యే చోట ఆగమనం ఉపయోగకరం. ప్రత్యేక పద్ధతిలో నిర్మించిన గణిత వస్తువులపై మాత్రమే ఆగమనం పనిచేస్తుంది: సమితిలోని ప్రతి మూలకం ప్రాథమికమైనదో, లేదా ప్రాథమిక మూలకాల నుంచి నిర్మితమైనదో అయితే దానిపై ఆగమనం చేయవచ్చు. మరింత ఖచ్చితంగా:
\end{explain}

\begin{defn}
సమితి $S$ను ప్రమేయం $f$ కింద \emph{సంవృతం} అంటాం, $x \in S$ అయితే $f(x) \in S$ అయినప్పుడు, అప్పుడే.
\end{defn}

\begin{explain}
ఉదాహరణకు, సహజ సంఖ్యలు ($\mathbb{N}$) ఉత్తరవర్తి ప్రమేయం $f(x) = x+1$ కింద సంవృతమైనవి. $x$ సహజ సంఖ్య అయితే $x+1$ కూడా సహజ సంఖ్యే. అయితే పూర్వవర్తి ప్రమేయం $f(x) = x-1$ కింద సహజ సంఖ్యలు సంవృతమైనవి కావు: $0$ సహజ సంఖ్య, కాని $-1$ కాదు.
\end{explain}

\begin{defn}
ప్రమేయం $f(x)$ విషయంలో, $f$ నిర్వచన సమితిలోని ఏ వస్తువు $a$కైనా $P(a)$ నుంచి $P(f(a))$ వస్తే, లక్షణం $P$ ఆ ప్రమేయం కింద \emph{నిలుపబడుతుంది} అంటారు ($a$కు $P$ ఉంటే $f(a)$కు కూడా ఉంటుంది).
\end{defn}

\begin{explain}
“సరి సంఖ్య కావడం” అనే లక్షణం $f(x) = x+2$ కింద నిలుపబడుతుంది; ఎందుకంటే $x$ సరి సంఖ్య అయితే $x+2$ కూడా సరి సంఖ్య. అదే లక్షణం ఉత్తరవర్తి ప్రమేయం కింద నిలుపబడదు; $x$ సరి సంఖ్య అయితే $x+1$ బేసి సంఖ్య.

ప్రతి సహజ సంఖ్యను 0 నుంచి ఉత్తరవర్తి ప్రమేయాన్ని పరిమిత సంఖ్యలో వర్తింపజేసి పొందవచ్చు. అందువల్ల సహజ సంఖ్యలపై ఆగమనం పనిచేస్తుంది. ఆగమనం చేయగల ఏ సమితికైనా ఈ నిర్మాణ లక్షణం ముఖ్యం.
\end{explain}

\begin{defn}[$\Nat$పై ఆగమనం]
0కు లక్షణం $P$ ఉండి, ఉత్తరవర్తి ప్రమేయం కింద $P$ నిలుపబడితే, ప్రతి సహజ సంఖ్యకూ $P$ ఉంటుంది.
\end{defn}

\begin{ex} సున్నా నుంచి $n$ వరకు సహజ సంఖ్యల మొత్తం $\frac{n(n+1)}{2}$. \\

\textbf{ఆధార దశ.} సున్నా వరకు సహజ సంఖ్యల మొత్తం $0=\frac{0\cdot 1}{2}$.\\

\textbf{ఆగమన దశ.} $k$కు లక్షణం నిజమని అనుకుందాం. అది $k+1$కూ నిజమని చూపుతాం. మరో మాటలో $P(k)$ నిజమని, అంటే
\[ 0 + 1 + \ldots + k = \frac{k(k+1)}{2} \]
అని అనుకుందాం. ఈ పరికల్పనను ఉపయోగించి $P(k+1)$, అంటే
\[ 0 + 1 + \ldots + k + (k+1) = \frac{(k+1)(k+2)}{2} \]
అని చూపుతాం; ఇందుకు $P(k)$ అనే పరికల్పనను ఉపయోగిస్తాం.

\begin{align*}
0 + 1 + \ldots + k &= \frac{k(k+1)}{2} \tag{పరికల్పన} \\
0 + 1 + \ldots + k + (k+1) &= \frac{k(k+1)}{2} + (k+1)
\tag{రెండు వైపులా $k+1$ కలిపాం} \\
&= \frac{k(k+1) + 2(k+1)}{2} \\
&= \frac{k^2 + k + 2k +2}{2} \\
&=\frac{(k+1)(k+2)}{2}
\end{align*}
కాబట్టి ఆగమన దశ పూర్తయింది; వాదనను నిరూపించాం.\sourcecorrection{OLTESFRINDINT-001}{మూలంలోని మూడో సమానతలో గుణకారాన్ని తప్పుగా విస్తరించారు: మొదటి పదం రెండింతల $k$ కాదు, $k$ యొక్క వర్గం. అదే సమానతలో సరి చేసిన రూపం తదుపరి $(k+1)(k+2)$ గుణకారానికి సరిపోతుంది.}
\end{ex}

\begin{explain}
\tetoken{సూత్రాల}{!!{formula}s} విషయంలో ప్రాథమిక మూలకాలు పరమాణు \tetoken{సూత్రాలు}{!!{formula}s}. సరళమైన \tetoken{సూత్రాలను}{!!{formula}s} సంచాలకాలతో కలిపి మరింత సంక్లిష్ట \tetoken{సూత్రాలను}{!!{formula}s} పొందుతాం. ప్రతి సంచాలకాన్ని, రెండు \tetoken{సూత్రాలను}{!!{formula}s} తీసుకుని వాటిని సంబంధిత సంచాలకంతో కలిపిన కొత్త \tetoken{సూత్రాన్ని}{!!{formula}} ఇచ్చే ప్రమేయంగా చూడవచ్చు. ఉదాహరణకు, $!A$, $!B$ అనే రెండు \tetoken{సూత్రాలను}{!!{formula}s} సంయోగంగా కలిపే ప్రమేయాన్ని $\mathcal E _\land (!A, !B) = !A \land !B$ అని రాయవచ్చు. వీటిని \emph{సూత్ర నిర్మాణ ప్రమేయాలు} అనవచ్చు. ఈ ప్రమేయాల కింద \tetoken{సూత్రాల}{!!{formula}s} సమితి సంవృతమైనది: $!A$, $!B$ \tetoken{సూత్రాలు}{!!{formula}s} అయితే, $\lnot !A$, $!A \land !B$ మొదలైనవి కూడా సూత్రాలే.
\sourcecorrection{OLTESFRINDINT-002}{ఆంగ్ల మూలంలోని గణిత సూత్ర ఉదాహరణల్లో పునరావృతమైన అసంపూర్ణ గుర్తును, పక్కనున్న $!B$కు అనుగుణంగా సూత్ర చరరాశి $!A$గా సరిచేశాం.}
\end{explain}

\begin{defn}[\tetoken{సూత్రాలపై}{!!{formula}s} ఆగమనం]
ప్రతి పరమాణు \tetoken{సూత్రానికి}{!!{formula}} లక్షణం $P$ ఉండి, \tetoken{సూత్ర}{!!{formula}}-నిర్మాణ ప్రమేయాల కింద $P$ నిలుపబడితే, ప్రతి \tetoken{సూత్రానికీ}{!!{formula}} లక్షణం $P$ ఉంటుంది.
\end{defn}

\begin{explain}
ఈ నిర్వచనం \tetoken{సూత్రాలపై}{!!{formula}s} ఆగమన నిరూపణలకు ఒక “విధానాన్ని” సూచిస్తుంది. ప్రతి \tetoken{సూత్రానికి}{!!{formula}} లక్షణం $P$ ఉందని చూపమంటే, కింది దశలను అనుసరించండి:\\

\textbf{ఆధార దశ.} $!A$ ఒక పరమాణు \tetoken{సూత్రం}{!!{formula}} అనుకుందాం. [\ldots] అందువల్ల $!A$కు లక్షణం $P$ ఉంది.

\textbf{ఆగమన దశ.} $!A$, $!B$ రెండూ లక్షణం $P$ గల \tetoken{సూత్రాలు}{!!{formula}s} అనుకుందాం.

$\lnot$ సందర్భం: [\ldots] అందువల్ల $\lnot !A$కు లక్షణం $P$ ఉంది.

$\land$ సందర్భం: [\ldots] అందువల్ల $!A \land !B$కు లక్షణం $P$ ఉంది.

$\lor$ సందర్భం: [\ldots] అందువల్ల $!A \lor !B$కు లక్షణం $P$ ఉంది.

$\lif$ సందర్భం: [\ldots] అందువల్ల $!A \lif !B$కు లక్షణం $P$ ఉంది.

$\liff$ సందర్భం: [\ldots] అందువల్ల $!A \liff !B$కు లక్షణం $P$ ఉంది.

$\lforall[]$ సందర్భం: [\ldots] అందువల్ల $\lforall[x] !A$కు లక్షణం $P$ ఉంది.

$\lexists[]$ సందర్భం: [\ldots] అందువల్ల $\lexists[x] !A$కు లక్షణం $P$ ఉంది. \\

అందువల్ల ప్రతి \tetoken{సూత్రానికి}{!!{formula}} లక్షణం $P$ ఉంది.
\end{explain}

\end{document}
